\section{Retained Indexing Experiment and Scenario Illustration}\label{app:scope-workflow-boundary}
The retained model-assisted experiment compares flat checklist R, structured checklist T and raw \scope{} traversal V. There are 512 roots per arm and three query judgments per root (1,536 judgments). The endpoint is the mean of each root's three binary joint outcomes, requiring a correct deployment decision and exact configuration counts. R has 624 joint successes (40.63\%); V has 494 (32.16\%). T has 693 observed successes and three missing queries belonging to one root, giving the retained 45.12--45.31\% missing-outcome range. V is worse on this joint endpoint; this is retained negative evidence.

Decision-correct but count-incorrect judgments number 537 (R), 476 (T) and 695 (V), out of 1,536 query opportunities. Observed queries have no schema/evidence-ID failures. Binary-decision successes are 1,161 (R), 1,169 (T, plus three missing) and 1,189 (V). The 1.11--1.30 percentage-point V--T binary difference is a secondary post-outcome decomposition, not proof of a better workflow. Counts are query judgments, not independent roots; percentages average three judgments per root. We do not infer causality from the error decomposition.

The separate 24-scenario exercise uses author-constructed source-reuse cases and author labels: seven Compatible, twelve Incompatible and five Unresolved. The original rule-application outputs and three-category confusion matrices remain in \path{experiments/reviewer-audit/scenario_confusions.json}. This archived illustration supplies neither participant data nor an accuracy advantage over ordinary expert reading.
